\documentclass[10pt,a4paper]{amsart}
\usepackage[margin=19mm]{geometry}
\usepackage{amsmath,amssymb,mathtools}
\usepackage[scr=rsfs,cal=euler]{mathalfa}
\usepackage{xfrac}
\usepackage[unicode,hidelinks,bookmarks=false]{hyperref}
\newcommand{\ie}{i.e.\ }
\newcommand{\cf}{cf.\ }
\newcommand{\resp}{respectively\ }
\newcommand{\Subsection}{\subsection}
\providecommand{\nobreakdash}{\nobreak\hbox{-}}
\setlength{\emergencystretch}{2em}
\multlinegap0pt
\usepackage{tikz}
\usetikzlibrary{cd}
\usepackage{enumitem}
\newtheorem{lemma}{Lemma}
\title[pullback lemma]{Bounded cohomology, quotient extensions, and hierarchical hyperbolicity: the pullback lemma for quotient extensions}
\author{Francesco Fournier-Facio, Giorgio Mangioni, Alessandro Sisto}
\date{Source: arXiv:2505.20462v4 (CC BY 4.0)}
\begin{document}
\maketitle

\noindent\emph{Source and licence.} Bounded cohomology, quotient extensions, and hierarchical hyperbolicity, by Francesco Fournier-Facio, Giorgio Mangioni, Alessandro Sisto. Version arXiv:2505.20462v4
(CC BY 4.0), distributed by arXiv under the Creative Commons Attribution 4.0 International licence,
http://creativecommons.org/licenses/by/4.0/, as recorded in the arXiv licence metadata for this exact
version and shown on the abstract page https://arxiv.org/abs/2505.20462v4. Full licence notice:
\texttt{licenses/arxiv-2505.20462-CC-BY-4.0.txt}.\par\smallskip

\noindent\textit{Editorial scope.} Counted unit: the article's pullback lemma for quotient extensions (source0.tex lines 335--338, source label \texttt{qce pullback}): with the notation of the displayed quotient-extension diagram, the pullback map on second cohomology sends $[\overline{E}]$ to $[E]$, so in particular a bounded quotient extension pulls back to a bounded extension. The article's complete original proof of it is delivered with it (source0.tex lines 340--368, one \texttt{proof} environment of 29 lines, adjacent to the statement, with no gap and no deferral). No mathematics is written, repaired or re-derived: the statement and the proof are the article's own lines, in the article's own notation, and no retained line is substituted, re-flowed or omitted, so nothing is removed and the package carries no omission record. Retained uncounted as the source-local context the proof needs: the displayed quotient-extension diagram with its label (lines 307--325, source label \texttt{diagram qce}), which the statement invokes and which is delivered together with the article's own diagram code. Every cross-reference of every retained line resolves inside this document, and no retained line contains a citation, so the wrapper carries no bibliography and no bibliography entry. No retained line contains a figure environment or an image-inclusion command; the retained lines do contain the article's own commutative-diagram code, which the wrapper typesets with the standard \texttt{tikz} and \texttt{tikz-cd} packages, and the package declares no asset dependency.

Editorial formatting only, in addition: an AMS \texttt{amsart} preamble that restates the article's own declarations and macros used by the retained lines, namely the \texttt{lemma} environment and the standard packages the article itself loads for its diagrams and lists. The retained environment prints in this document as Lemma 1, whereas the article prints the same item with the numbering of its own full document; the one displayed numbered relation of the retained lines prints as (1), whereas the article numbers its equations per section.

The complete native source of the article as distributed in the arXiv e-print for this exact version is bundled unchanged as \texttt{sources/178803/source0.tex}: it is byte identical to the e-print file \texttt{main.tex} except that the pipeline's mechanical concatenation prepends one header comment line naming it. The article uses the AMS \texttt{amsart} class; no publisher class or style file is shipped in the e-print and none is redistributed or used.
\begin{equation}
\label{diagram qce}
    % https://q.uiver.app/#q=WzAsMTAsWzAsMCwiMSJdLFsxLDAsIloiXSxbMiwwLCJFIl0sWzMsMCwiRyJdLFs0LDAsIjEiXSxbMCwxLCIxIl0sWzEsMSwiWiJdLFsyLDEsIlxcb3ZlcmxpbmV7RX0iXSxbMywxLCJcXG92ZXJsaW5le0d9Il0sWzQsMSwiMSJdLFswLDFdLFsxLDJdLFsyLDMsIlxccGkiXSxbMyw0XSxbNSw2XSxbNiw3XSxbNyw4XSxbOCw5XSxbMSw2LCI9IiwyXSxbMiw3LCJwIiwyXSxbMyw4LCJwIiwyXV0=
\begin{tikzcd}
	1 & K & E & G & 1 \\
	1 & K & {\overline{E}} & {\overline{G}} & 1
	\arrow[from=1-1, to=1-2]
	\arrow[from=1-2, to=1-3]
	\arrow["{=}"', from=1-2, to=2-2]
	\arrow["\pi", from=1-3, to=1-4]
	\arrow["p"', from=1-3, to=2-3]
	\arrow[from=1-4, to=1-5]
	\arrow["p"', from=1-4, to=2-4]
	\arrow[from=2-1, to=2-2]
	\arrow[from=2-2, to=2-3]
	\arrow[from=2-3, to=2-4]
	\arrow[from=2-4, to=2-5]
\end{tikzcd}
\end{equation}

\begin{lemma}
\label{qce pullback}
    With the notation of Diagram \eqref{diagram qce}, the pullback $p^* \colon H^2(\overline{G}; K) \to H^2(G; K)$ sends $[\overline{E}]$ to $[E]$. In particular, if $[\overline{E}]$ is bounded, then $[E]$ is bounded.
\end{lemma}

\begin{proof}
    We choose a normalised section $\overline \sigma \colon \overline G \to \overline E$. We also choose an injective map $\tau \colon \overline G \to E$ such that $p \tau = \overline \sigma$ and $\tau(1_{\overline{G}}) = 1_E$. This way Diagram \eqref{diagram qce} is enriched as follows:
% https://q.uiver.app/#q=WzAsMTAsWzAsMCwiMSJdLFsxLDAsIloiXSxbMiwwLCJFIl0sWzMsMCwiRyJdLFs0LDAsIjEiXSxbMCwxLCIxIl0sWzEsMSwiWiJdLFsyLDEsIlxcb3ZlcmxpbmV7RX0iXSxbMywxLCJcXG92ZXJsaW5le0d9Il0sWzQsMSwiMSJdLFswLDFdLFsxLDJdLFsyLDMsIlxccGkiXSxbMyw0XSxbNSw2XSxbNiw3XSxbNyw4LCJcXG92ZXJsaW5lIFxccGkiLDAseyJvZmZzZXQiOi0xfV0sWzgsOV0sWzEsNiwiPSIsMl0sWzIsNywicCIsMl0sWzMsOCwicCJdLFs4LDcsIlxcb3ZlcmxpbmVcXHNpZ21hIiwwLHsib2Zmc2V0IjotMX1dLFs4LDIsIlxcdGF1IiwyXV0=
\[\begin{tikzcd}
	1 & K & E & G & 1 \\
	1 & K & {\overline{E}} & {\overline{G}} & 1
	\arrow[from=1-1, to=1-2]
	\arrow[from=1-2, to=1-3]
	\arrow["{=}"', from=1-2, to=2-2]
	\arrow["\pi", from=1-3, to=1-4]
	\arrow["p"', from=1-3, to=2-3]
	\arrow[from=1-4, to=1-5]
	\arrow["p", from=1-4, to=2-4]
	\arrow[from=2-1, to=2-2]
	\arrow[from=2-2, to=2-3]
	\arrow["{\overline \pi}", shift left, from=2-3, to=2-4]
	\arrow["\tau"', from=2-4, to=1-3]
	\arrow["{\overline\sigma}", shift left, from=2-4, to=2-3]
	\arrow[from=2-4, to=2-5]
\end{tikzcd}\]
    Now $\pi \tau$ is a section for $p \colon G \to \overline G$, so every element of $G$ can be written uniquely as $g = \pi(n) \cdot \pi \tau(p(g))$ for some $n\in N$. We define $\sigma(g) = n \cdot \tau(p(g))$, so $\sigma \colon G \to E$.
    First, note that $\sigma$ is indeed a normalised section: $\sigma(1_G) = 1_E$, and if $g = \pi(n) \pi \tau (p(g))$, then
    \[\pi\sigma(g) = \pi(n \cdot \tau(p(g))) = \pi(n) \cdot \pi \tau(p(g)) = g.\]
    Secondly, we claim that $p \sigma = \overline \sigma p$, indeed
    \[p\sigma(g) = p(n \cdot \tau(p(g))) = p\tau(p(g)) = \overline \sigma(p(g)),\]
    which implies that the projection $p \colon E \to \overline E$ sends
    \[\sigma(g)\sigma(h)\sigma(gh)^{-1} \mapsto \overline \sigma(p(g)) \overline \sigma(p(h)) \overline \sigma(p(gh))^{-1}.\]
    This shows that the cocycle defined by $\sigma$ is indeed the pullback of the cocycle defined by $\overline \sigma$, and concludes the proof.
\end{proof}

\end{document}
